%% Paper B, REVISION 1 for C&I (submission 2067810, decision of 2026-11-06).
%%
%% Working copy of ../../paper_b_ci.tex as submitted on 2026-10-21 (that file and
%% ../../submission/ stay frozen; hashes in ../SUBMITTED_SNAPSHOT.sha256). Every change
%% made here is recorded in ../CHANGELOG.md with the reviewer items it answers, and
%% ../scripts/build_revision.py renders it and a changes-marked PDF into ../build/.
%% Paths below are relative to ci_journal/, where the build compiles it.
%%
%% DRAFT of 2026-11-06 (decision D-D): prepared by the assistant from the approved
%% proposal (P-01..P-24, PROPOSAL_2026-11-06.md) for the author to edit. Text marked
%% [R: ...] is filled from the results of cohort B2; [REF: ...] marks literature that the
%% reference audit (WP8) supplies with verified DOIs. No number of the submitted version
%% is reused.
%%
%% This file is not compiled for the journal (it accepts Word only):
%% scripts/build_paper_b_ci_docx.py writes it into the journal template.
\documentclass[11pt,a4paper]{article}
\usepackage{amssymb}
\usepackage{amsmath}
\usepackage{graphicx}
\usepackage{booktabs}
\usepackage{tabularx}
\usepackage{array}
\usepackage{enumitem}
\usepackage{url}
\usepackage[hidelinks]{hyperref}
\setlength{\emergencystretch}{3em}
\setlist{nosep,leftmargin=*}

% The journal reviews masked (double-blind): the author, the affiliation,
% the acknowledgments, the repository address and the archive DOI are on
% the separate title page only.

\title{When Does SHAP Outperform LIME? Model- and Configuration-Dependent
Results from a Paired Tabular Benchmark}
\author{}
\date{}

\begin{document}
\maketitle

\begin{abstract}
\noindent
\input{revision_1/manuscript/abstract_rev1.tex}
\end{abstract}

% Keywords: revision_1/manuscript/keywords_rev1.tex; the builder puts them in
% alphabetical order with a capital first letter, as the journal asks.

% -----------------------------------------------------------------------
\section{Introduction}
\label{sec:intro}
% -----------------------------------------------------------------------

A team that deploys a tabular classifier and must explain its predictions
faces a concrete decision: which post-hoc explainer to use, for which model,
under which latency budget. The candidates differ in the kind of explanation
they return. LIME~\cite{ribeiro2016why} and SHAP~\cite{lundberg2017unified}
return feature attributions, Anchors~\cite{ribeiro2018anchors} returns a
rule and DiCE~\cite{mothilal2020dice} a counterfactual. They also differ in
how faithfully the explanation tracks the model, how stable it is under
small changes of the input, and how much computation it needs. Theory does
not settle the choice across model families and
configurations~\cite{retzlaff2024posthoc,ali2023trustworthy,longo2024manifesto}, and the
empirical guidance is weaker than the popularity of the methods suggests:
evaluations use different metric families, and the ranking of explainers
often depends as much on the metric and the configuration as on the
methods~\cite{nauta2023anecdotal,canha2025benchmark,pawlicki2024metrics}.

Two gaps follow. Comparisons in which all explainers see the same model,
the same instances and the same seed, so that a difference can be
attributed to the explainer, are scarce. Results are also usually reported
for one configuration, although the hyperparameters that build LIME's
neighborhood or bound KernelSHAP's sampling can change the outcome. Cost,
moreover, is mostly reported as wall-clock time on unstated hardware, which
another team cannot transfer to its own deployment.

This paper addresses these gaps with three contributions:
\begin{enumerate}
\item \textbf{A paired benchmark of four explainers} (SHAP, LIME, Anchors and
DiCE) on five model families, with every run on one recorded host, and cost
reported as CPU time and as the number of model evaluations, a measure that
does not depend on the hardware (Sections~\ref{sec:main}
and~\ref{sec:cost}).
\item \textbf{Configuration-aware results}: a sweep over the LIME kernel
width and sample size and over the KernelSHAP budget and background, and
controls that test whether the stability measurement itself is valid
(Sections~\ref{sec:sweep} and~\ref{sec:validity_controls}).
\item \textbf{Bounded recommendations}: every claim is labelled as
confirmatory or exploratory and stated with the models, datasets and
configurations under which it holds (Section~\ref{sec:discussion}).
\end{enumerate}
The paper makes ranking claims under a stated protocol. It does not claim
that an explanation is correct for a task, and it does not test whether
people decide better with any explainer.

% -----------------------------------------------------------------------
\section{Background and Related Work}
\label{sec:background}
% -----------------------------------------------------------------------

LIME fits an interpretable surrogate $g$ to the model $f$ on perturbed
samples around an instance $x$, weighted by a proximity kernel $\pi_x$, as
in Eq.~(\ref{eq:lime}):
\begin{equation}
\xi(x) = \arg\min_{g \in G} \mathcal{L}(f, g, \pi_x) + \Omega(g),
\label{eq:lime}
\end{equation}
where $G$ is the class of surrogates (sparse linear models here),
$\mathcal{L}$ measures how poorly $g$ reproduces $f$ on the weighted
samples and $\Omega(g)$ penalizes surrogate complexity~\cite{ribeiro2016why}.
The kernel width, the number of samples and the feature budget shape the
neighborhood, and the sampling makes the explanation stochastic. SHAP
assigns additive attributions $f(x) = \phi_0 + \sum_{i=1}^{M} \phi_i$ that
satisfy local accuracy, missingness and consistency, where $\phi_0$ is the
expected output over a background sample and $M$ the number of
features~\cite{lundberg2017unified}. KernelSHAP estimates them by a
weighted regression over sampled feature coalitions and applies to any
model; TreeSHAP computes them from the structure of tree ensembles. Anchors
searches for a rule that fixes the prediction with a target
precision~\cite{ribeiro2018anchors}; DiCE searches for a counterfactual
instance with the opposite prediction~\cite{mothilal2020dice}.

Quantitative evaluation relies on proxy metrics: faithfulness under feature
removal~\cite{bhatt2020evaluating,samek2017evaluating,hooker2019benchmark},
robustness to small input changes~\cite{alvarezmelis2018robustness} and
sparsity, standardized in toolkits and
benchmarks~\cite{hedstrom2023quantus,agarwal2022openxai}. Removal-based
scores depend on the removal scheme~\cite{rong2022consistent,zheng2025ffidelity}, attributions
can pass perturbation tests and fail sanity checks~\cite{adebayo2018sanity}, LIME and SHAP can be manipulated~\cite{slack2020fooling},
and Shapley attributions are hard to read as importance when features are
dependent~\cite{kumar2020problems}. [REF: WP8 adds a paragraph on empirical
comparisons of SHAP and LIME and on explanation benchmarks published from
2024 on (verified DOIs), and positions this study against them: what they
compared, under which pairing and configuration, and how cost was
reported.] [REF: work on the sensitivity of LIME to its kernel width and
sample size, and on the validity of perturbation-based stability measures
for encoded features, which motivates the controls of
Section~\ref{sec:validity_controls}.]

% -----------------------------------------------------------------------
\section{Materials and Methods}
\label{sec:methods}
% -----------------------------------------------------------------------

\subsection{Design and Research Questions}
\label{sec:rqs}

All results come from one experimental cohort run for this study under the
protocol below (Table~\ref{tab:design}). An earlier analysis of the
SHAP--LIME question used runs collected on several hosts; an audit found that
in most of its cells the evaluation instances overlapped the training
partition, that one model could not be recovered and that its timings mixed
hosts, so no number from it is used here.

For each endpoint $Y$ and each instance, the paired difference is
$\Delta_X = Y_{\mathrm{SHAP}} - Y_X$ for $X \in \{\mathrm{LIME},
\mathrm{Anchors}, \mathrm{DiCE}\}$. Three directional hypotheses are tested,
and sparsity is reported descriptively:
\begin{itemize}
\item \textbf{H1 (stability).} SHAP is more stable than $X$:
$\Delta_X(S) > 0$.
\item \textbf{H2 (cost).} SHAP is more expensive than $X$:
$\Delta_X(C) > 0$, in CPU time and in model evaluations.
\item \textbf{H3 (faithfulness).} SHAP is more faithful than $X$:
$\Delta_X(F) > 0$ and $\Delta_X(G) > 0$.
\item \textbf{Sparsity (descriptive).} LIME's sparsity is fixed by its
feature budget and Anchors returns few features by construction, so a test
of sparsity would test a configuration, not a method.
\end{itemize}

\begin{table}[htb]
\centering
\caption{Design of the benchmark. The primary benchmark is confirmatory; the
other parts are exploratory}
\label{tab:design}
\begin{tabularx}{\linewidth}{p{3.1cm} X}
\toprule
\textbf{Factor} & \textbf{Levels} \\
\midrule
Datasets & Adult (primary, 108 encoded features); Breast Cancer (30), German
Credit (61), Diabetes (8), Spambase (57), Phishing (68), Connect-4 (126) \\
Models & logistic regression, random forest, XGBoost, SVM (RBF), MLP \\
Explainers & SHAP, LIME (full vector; 10-feature variant), Anchors, DiCE \\
Seeds & Adult: 5; other datasets: 3 \\
Instances & up to 25 per stratum (TP, TN, FP, FN) of the test split, per
model and seed; stability on 3 per stratum \\
Paired cells & Adult: 25 (model $\times$ seed); other datasets: 90 \\
Timing subset & 5 instances per stratum, 3 repetitions, one host \\
Sweep and controls & Section~\ref{sec:sweep} and
Section~\ref{sec:validity_controls} \\
\bottomrule
\end{tabularx}
\end{table}

\subsection{Data and Models}
\label{sec:data}

Adult~\cite{kohavi1996scaling,becker1996adult} is the primary dataset. Six
further binary datasets were chosen so that feature dimension crosses
feature type: Breast Cancer~\cite{wolberg1993breast} and German
Credit~\cite{hofmann1994statlog} from the UCI repository~\cite{dua2019uci}, and Diabetes, Spambase,
Phishing and Connect-4 [REF: dataset citations] in a $2 \times 2$ of
numeric or categorical features and lower or higher dimension. For
Connect-4, a win is the positive class and the two other outcomes the
negative class.

Each dataset has one fixed stratified split (80/20, seed 42), shared by
preprocessing, model training and every seed. Numeric features are scaled
and categorical features one-hot encoded, both fitted on the training
split only. Training is capped at 20,000 rows (a stratified subsample) so
that the SVM fits in reasonable time. The five
models~\cite{breiman2001random,chen2016xgboost} are trained once per dataset
with fixed hyperparameters (supplementary file) and frozen; every run checks
the model file against its recorded hash. A seed therefore changes the
evaluation instances and the explainers' randomness, never the model.
Evaluation instances are drawn from the test split only, up to 25 from each
confusion-matrix stratum; smaller strata are taken whole, and the realized
counts are reported. Within a cell, all explainers see the same instances.
In small strata, different seeds can draw the same instance; the overlap
is measured and reported [R: overlap].

\subsection{Explainers}
\label{sec:explainers}

All explainers receive the same encoded features and explain the predicted
probability of the positive class:
\begin{itemize}
\item \textbf{SHAP:} TreeSHAP for the random forest and XGBoost;
KernelSHAP for the other models, with 50 background rows and a budget of
500 coalitions.
\item \textbf{LIME:} 1,000 samples and the library's default kernel width,
$0.75\sqrt{M}$ ($\approx 7.79$ on Adult). The primary contrast uses the full
attribution vector; a variant with the 10-feature budget shows the effect
of the budget.
\item \textbf{Anchors:} precision target 0.95; the attribution vector marks
the features of the anchor with 1 and the others with 0.
\item \textbf{DiCE:} one counterfactual of the opposite class (random
search); the attribution of feature $j$ is $|x_j - x'_j|$.
\end{itemize}
These settings are the libraries' defaults where defaults exist, so that a
practitioner's first run is what is measured. KernelSHAP's budget is set
explicitly because the default ($2M + 2048$ coalitions) made the SVM the
cost driver of the whole design; Section~\ref{sec:sweep} measures the effect
of the default against the explicit budget. Anchors and DiCE return
explanations of different kinds, so their vectors measure which features the
explanation uses, not signed contributions; the metrics below are
comparable across the four methods in that sense.

\subsection{Metrics}
\label{sec:metrics}

Let $f(x)$ be the predicted probability of the positive class and
$e(x) \in \mathbb{R}^{M}$ the attribution vector; masking feature $j$
replaces it with its training mean $\bar{x}_j$. The metrics follow common
definitions~\cite{alvarezmelis2018robustness,bhatt2020evaluating,samek2017evaluating}:
\begin{enumerate}
\item \textbf{Stability} $S$: the mean cosine similarity over all pairs of
$T = 15$ explanations of perturbed copies $x + \varepsilon^{(a)}$,
$\varepsilon^{(a)} \sim \mathcal{N}(0, \sigma^2 I)$, $\sigma = 0.1$,
Eq.~(\ref{eq:stability}):
\begin{equation}
S = \frac{2}{T(T-1)}\sum_{1 \le a<b \le T}\cos\!\left(e^{(a)},e^{(b)}\right).
\label{eq:stability}
\end{equation}
The cosine with an all-zero vector is set to 0; such vectors are counted
and $S$ is also reported without them.
\item \textbf{Fidelity} $F$: the Pearson correlation, across the $M$
features, between $|e_j(x)|$ and the prediction change
$|f(x) - f(\tilde{x}^{(j)})|$ when feature $j$ alone is masked,
Eq.~(\ref{eq:fidelity}):
\begin{equation}
F = \rho\!\left(\bigl\{|e_j(x)|\bigr\}_{j=1}^{M},\;
\bigl\{|f(x) - f(\tilde{x}^{(j)})|\bigr\}_{j=1}^{M}\right).
\label{eq:fidelity}
\end{equation}
Absolute values are used because the masked change is unsigned. When
either vector is constant, $F$ is undefined and set to 0; these cases are
counted.
\item \textbf{Faithfulness gap} $G$: the prediction change when the $k = 5$
features with the largest $|e_j(x)|$ are masked together.
\item \textbf{Sparsity} $P$: the share of features with
$|e_j(x)| > 10^{-4}$ (lower is sparser).
\item \textbf{Cost} $C$: CPU time per explanation and the number of rows
the explainer passes to the model (Section~\ref{sec:cost_method}).
\end{enumerate}

\subsection{Statistical Analysis}
\label{sec:inference}

Instances are nested in seeds and seeds in models, so instance-level
differences are not independent. For each endpoint and contrast, the
instance-level differences $\Delta_{gsi}$ of model $g$, seed $s$ and
instance $i$ are fitted by a linear mixed model
$\Delta_{gsi} = \mu + u_g + v_{gs} + \epsilon_{gsi}$, with random intercepts
for the model and for the seed within the model~\cite{bates2015lme4}; $\mu$
is the effect, reported with its 95\% confidence interval. Because five
models are few clusters, the interval is checked by a cluster bootstrap
that resamples models~\cite{field2007bootstrap}. Next to the standardized
paired difference $d_z$ of the cell means~\cite{lakens2013calculating}, the
Hodges--Lehmann estimate~\cite{hodges1963estimates} and the matched-pairs
rank-biserial correlation~\cite{kerby2014simple}, the effect size of the Wilcoxon signed-rank test~\cite{wilcoxon1945individual}, are reported. The
confirmatory family is three endpoints (stability, faithfulness as $F$ and
$G$, cost) times three contrasts; $p$-values are Holm-corrected over the
twelve tests~\cite{holm1979simple}. The four explainers are ranked on cell
means by the Friedman test with the Nemenyi critical
difference~\cite{demsar2006statistical}. The sweep, the controls and the six
further datasets are exploratory and are reported without tests.

\subsection{Cost Measurement and Computing Environment}
\label{sec:cost_method}

All runs used one laptop (AMD Ryzen 5 4600H, 6 cores, 31.4 GB RAM,
Windows 11), Python 3.11 with scikit-learn 1.7.1, XGBoost 3.1.2, SHAP
0.50.0, LIME 0.2.0.1, alibi 0.9.6 (Anchors) and dice-ml 0.12. Quality was
computed in parallel worker processes; quality does not depend on speed.
Cost was measured in a separate pass on a subset (Table~\ref{tab:design}):
one process with one thread, pinned to one core, on mains power; the
explainers interleaved per instance, after one warm-up explanation each;
three repetitions, of which the median is used; and measurements paused
while any other process used more than 10\% of a core or the machine was
in use. A wrapper counts every row the explainer passes to the model.
Rows evaluated, unlike time, transfer to other hardware: multiplied by the
cost of one model evaluation they give the explainer's cost on any host.
For logistic regression and the MLP, whose cost per row is known, this is
also given in floating-point operations.

% -----------------------------------------------------------------------
\section{Results}
\label{sec:results}
% -----------------------------------------------------------------------

\subsection{Primary Benchmark on Adult}
\label{sec:main}

Table~\ref{tab:paired_main} reports the effects on the 25 Adult cells.
[R: per contrast, the direction and size of the stability and faithfulness
effects, and in how many cells the sign holds.] [R: Friedman ranking of
the four explainers per endpoint.] [R: share of all-zero vectors and of
undefined fidelity per explainer, and whether excluding them changes any
effect.] The distribution of the instance-level differences is shown in
the supplementary file.

\begin{table}[htb]
\centering
\caption{Primary benchmark (Adult, 25 cells). Per-method mean levels and
paired effects $\Delta_X = $ SHAP $- X$: mixed-model estimate with 95\% CI,
Holm-adjusted $p$, standardized paired difference $d_z$, Hodges--Lehmann
estimate (HL) and matched rank-biserial correlation $r_{rb}$. Cost: CPU
ms per explanation}
\label{tab:paired_main}
\begin{tabular}{llrrrrrrr}
\toprule
Endpoint & $X$ & SHAP & $X$ & $\Delta_X$ [95\% CI] & $p_{\mathrm{Holm}}$ & $d_z$ & HL & $r_{rb}$ \\
\midrule
Stability & LIME & [R] & [R] & [R] & [R] & [R] & [R] & [R] \\
 & Anchors & & [R] & [R] & [R] & [R] & [R] & [R] \\
 & DiCE & & [R] & [R] & [R] & [R] & [R] & [R] \\
Fidelity & LIME & [R] & [R] & [R] & [R] & [R] & [R] & [R] \\
 & Anchors & & [R] & [R] & [R] & [R] & [R] & [R] \\
 & DiCE & & [R] & [R] & [R] & [R] & [R] & [R] \\
Faithf.\ gap & LIME & [R] & [R] & [R] & [R] & [R] & [R] & [R] \\
 & Anchors & & [R] & [R] & [R] & [R] & [R] & [R] \\
 & DiCE & & [R] & [R] & [R] & [R] & [R] & [R] \\
Cost & LIME & [R] & [R] & [R] & [R] & [R] & [R] & [R] \\
 & Anchors & & [R] & [R] & [R] & [R] & [R] & [R] \\
 & DiCE & & [R] & [R] & [R] & [R] & [R] & [R] \\
\midrule
Sparsity & all & [R] & [R] & \multicolumn{5}{l}{descriptive} \\
\bottomrule
\end{tabular}
\end{table}

\subsection{Cost}
\label{sec:cost}

Table~\ref{tab:cost} gives the cost per explanation by model family.
[R: CPU-time ordering by model family; the share of model time in the
explainer's time.] [R: rows evaluated per explanation: LIME 1,000; KernelSHAP
about 25,000 (500 coalitions $\times$ 50 background rows); Anchors and DiCE as
measured.] Earlier runs suggested that TreeSHAP was faster than LIME on
XGBoost but slower on the random forest, two ensembles of different size.
Here both have 100 trees, and the comparison is repeated with matched depth and with cost
plotted against the total number of leaves [R: whether the reversal holds,
and whether it follows the size of the forest or the implementation].

\begin{table}[htb]
\centering
\caption{Cost per explanation on Adult (one host, one core): median CPU time
and rows evaluated by the model; theoretical complexity per explanation
($T$ trees, $L$ leaves, $D$ depth; $n_s$ samples or coalitions; $n_b$
background rows). FLOPs for logistic regression and the MLP in the text}
\label{tab:cost}
\begin{tabular}{lllrr}
\toprule
Explainer & Model & Complexity & CPU (ms) & Rows evaluated \\
\midrule
SHAP (Tree) & RF, XGB & $O(T L D^2)$, no model calls & [R] & 0 \\
SHAP (Kernel) & logreg, SVM, MLP & $O(n_s n_b)$ rows & [R] & [R] \\
LIME & all & $O(n_s)$ rows $+$ fit $O(n_s M^2)$ & [R] & [R] \\
Anchors & all & adaptive (precision target) & [R] & [R] \\
DiCE & all & iterative search & [R] & [R] \\
\bottomrule
\end{tabular}
\end{table}

\subsection{Configuration Sensitivity}
\label{sec:sweep}

The sweep varied, on Adult for every model and three seeds, LIME's kernel
width and number of samples, and KernelSHAP's budget and background size;
KernelSHAP was also run on the two tree ensembles, to separate the method
from the TreeSHAP implementation. [R: Figure~1 placeholder: Pareto fronts of
stability, fidelity and cost per model, one panel per model, computed on the
instances valid under every configuration.] [R: whether the sign of each
SHAP--LIME difference holds over the grid; the share of degenerate
explanations at small kernel widths; the effect of the library default
against the 500-coalition budget.]

\subsection{Validity of the Measurements}
\label{sec:validity_controls}

Five controls test whether the stability and fidelity measurements measure
the explainer: (a) explanations of the unperturbed instance repeated $T$
times, which isolates each explainer's own sampling variance; (b) noise
that keeps one-hot features valid (a category is resampled instead of adding
Gaussian noise to its indicators); (c) rank correlation and top-$k$ overlap
on the full vectors, besides the cosine; (d) the effect of excluding
all-zero vectors; and (e) the fidelity ordering under four masking schemes
(zero, mean, marginal and grouped replacement, where associated features
such as education and its numeric code are masked together) on every model.
[R: for each control, whether the SHAP--LIME stability and fidelity
orderings hold, and how much of LIME's instability is its own sampling
variance.]

\subsection{Other Datasets}
\label{sec:cross}

Table~\ref{tab:cross} reports the four explainers on the six further
datasets, now including SHAP stability. [R: whether the Adult orderings
hold across the $2 \times 2$ of dimension and feature type; where they
change.]

\begin{table}[htb]
\centering
\caption{Other datasets (exploratory): mean $\pm$ SD over 3 seeds, averaged
over the five models (per model in the supplementary file)}
\label{tab:cross}
\begin{tabular}{lcccccccc}
\toprule
 & \multicolumn{4}{c}{Stability} & \multicolumn{4}{c}{Fidelity} \\
Dataset & SHAP & LIME & Anch. & DiCE & SHAP & LIME & Anch. & DiCE \\
\midrule
Breast Cancer & [R] & [R] & [R] & [R] & [R] & [R] & [R] & [R] \\
German Credit & [R] & [R] & [R] & [R] & [R] & [R] & [R] & [R] \\
Diabetes & [R] & [R] & [R] & [R] & [R] & [R] & [R] & [R] \\
Spambase & [R] & [R] & [R] & [R] & [R] & [R] & [R] & [R] \\
Phishing & [R] & [R] & [R] & [R] & [R] & [R] & [R] & [R] \\
Connect-4 & [R] & [R] & [R] & [R] & [R] & [R] & [R] & [R] \\
\bottomrule
\end{tabular}
\end{table}

% -----------------------------------------------------------------------
\section{Discussion}
\label{sec:discussion}
% -----------------------------------------------------------------------

\subsection{What Holds, and Where}

Table~\ref{tab:holds} states each claim with its evidence and its scope.
Confirmatory claims rest on the primary benchmark and its corrected tests;
exploratory claims rest on the sweep, the controls and the further
datasets, and are reported as patterns, not as tested effects. [R: two or
three sentences on the main pattern, including the conditions under which an
ordering changes.]

\begin{table}[htb]
\centering
\caption{What holds, where, and under which configuration}
\label{tab:holds}
\begin{tabularx}{\linewidth}{p{3.0cm} p{2.0cm} X}
\toprule
\textbf{Claim} & \textbf{Status} & \textbf{Scope} \\
\midrule
SHAP more stable than $X$ & confirmatory & [R] \\
SHAP more faithful than $X$ & confirmatory & [R] \\
SHAP more expensive than $X$ & confirmatory & [R] \\
Orderings over the configuration grid & exploratory & [R] \\
Orderings on other datasets & exploratory & [R] \\
TreeSHAP vs LIME cost by ensemble & exploratory & [R] \\
Usefulness for people & not tested & no participant study \\
\bottomrule
\end{tabularx}
\end{table}

The recommendations are restricted to what was tested. [R: where SHAP is
preferable on faithfulness and stability and at what cost; where a cheaper
explainer suffices; which configuration choices a practitioner should check
before relying on either rule.]

\subsection{Relation to Earlier Work}

Earlier comparisons of SHAP and LIME, including a multi-metric benchmark on
Adult~\cite{herrera2026framework}, reported SHAP ahead on fidelity and
stability and LIME ahead on cost, for one configuration and with timings
from unrecorded hardware. This study differs in four respects: all four
explainers are paired on the same instances; every run uses one recorded
host and cost is also counted in model evaluations; the configuration is
varied and the measurement is tested; and the results extend to six further
datasets. [R: which earlier statements are confirmed, qualified or
reversed.]

\subsection{Limitations}
\label{sec:validity}

The endpoints compare explainers under a protocol; they do not establish
that an explanation is correct, which needs ground truth or
human-centered evaluation~\cite{haufe2026formalization,doshivelez2017rigorous,jacovi2020faithfulness},
and attribution methods can mark suppressor variables as important when
features are dependent~\cite{wilming2022scrutinizing,wilming2023theoretical}.
Anchors and DiCE vectors encode which features an explanation uses, not
signed contributions. Cost was measured on one host; times scale with the
hardware, rows evaluated do not. The study covers tabular binary
classification with five model families and seven datasets.

% -----------------------------------------------------------------------
\section{Conclusions}
\label{sec:conclusion}
% -----------------------------------------------------------------------

[R: two to three sentences that restate only results of
Table~\ref{tab:paired_main} and Table~\ref{tab:holds}: the confirmatory
orderings with effect sizes, the cost in rows evaluated, and the conditions
under which an ordering changes.] The choice of explainer depends on the
model architecture, the configuration and the latency budget, and a ranking
obtained under one configuration should not be transferred to another
without measuring it. Whether any of these explainers helps people decide
is a separate question that needs a study with participants.

% The acknowledgment is on the separate title page (masked review).

\section*{Competing Interests}

No competing financial interests or personal relationships that could have
appeared to influence the work reported in this paper are declared. No
external funding was received for this study.

\section*{Data and Code Availability}
\label{sec:artifacts}

The code, the frozen models with their hashes, every result row and the
analysis scripts are available in a public repository and a versioned
archive with a DOI, under the MIT License; their addresses are withheld for
masked review and are given on the title page. A supplementary file lists
the design, the hyperparameters, the per-dataset and per-model results and
the commands that regenerate every table from the result rows.

\begin{thebibliography}{99}

\bibitem{ribeiro2016why}
Ribeiro MT, Singh S, Guestrin C. ``Why should I trust you?'': Explaining the
predictions of any classifier. In: Proceedings of the 22nd ACM SIGKDD
International Conference on Knowledge Discovery and Data Mining; 2016.
p. 1135--1144. \url{https://doi.org/10.1145/2939672.2939778}

\bibitem{lundberg2017unified}
Lundberg SM, Lee SI. A unified approach to interpreting model predictions.
In: Advances in Neural Information Processing Systems 30; 2017. p. 4765--4774.

\bibitem{retzlaff2024posthoc}
Retzlaff CO, Angerschmid A, Saranti A, Schneeberger D, R\"{o}ttger R,
M\"{u}ller H, et al. Post-hoc vs ante-hoc explanations: xAI design guidelines
for data scientists. Cognitive Systems Research. 2024;86:101243.
\url{https://doi.org/10.1016/j.cogsys.2024.101243}

\bibitem{ali2023trustworthy}
Ali S, Abuhmed T, El-Sappagh S, Muhammad K, Alonso-Moral JM, Confalonieri R,
et al. Explainable Artificial Intelligence (XAI): What we know and what is
left to attain Trustworthy Artificial Intelligence. Information Fusion.
2023;99:101805. \url{https://doi.org/10.1016/j.inffus.2023.101805}

\bibitem{nauta2023anecdotal}
Nauta M, Trienes J, Pathak S, Nguyen E, Peters M, Schmitt Y, et al. From
anecdotal evidence to quantitative evaluation methods: A systematic review on
evaluating explainability. ACM Computing Surveys. 2023;55(13s):1--42.
\url{https://doi.org/10.1145/3583558}

\bibitem{canha2025benchmark}
Canha D, Kubler S, Fr\"{a}mling K, Fagherazzi G. A functionally-grounded
benchmark framework for XAI methods: Insights and foundations from a
systematic literature review. ACM Computing Surveys. 2025;57(12):1--40.
\url{https://doi.org/10.1145/3737445}

\bibitem{longo2024manifesto}
Longo L, Brcic M, Cabitza F, Choi J, Confalonieri R, Del Ser J, et al.
Explainable artificial intelligence (XAI) 2.0: A manifesto of open challenges
and interdisciplinary research directions. Information Fusion.
2024;106:102301. \url{https://doi.org/10.1016/j.inffus.2024.102301}

\bibitem{pawlicki2024metrics}
Pawlicki M, Pawlicka A, Uccello F, Szelest S, D'Antonio S, Kozik R, et al.
Evaluating the necessity of the multiple metrics for assessing explainable
AI: A critical examination. Neurocomputing. 2024;602:128282.
\url{https://doi.org/10.1016/j.neucom.2024.128282}

\bibitem{zheng2025ffidelity}
Zheng X, Shirani F, Chen Z, Lin C, Cheng W, Guo W, et al. F-Fidelity: A
robust framework for faithfulness evaluation of explainable AI. In: The 13th
International Conference on Learning Representations (ICLR); 2025.
\url{https://openreview.net/forum?id=X0r4BN50Dv}

\bibitem{herrera2026framework}
Herrera-V\'asquez J, Herrero-Uceda M. A framework for rigorous evaluation of
model-agnostic explainability methods: multi-metric statistical benchmarking,
operational protocol, and reproducibility. Revista de Investigaci\'on
Multidisciplinaria Iberoamericana (RIMI). 2026;(3).
\url{https://doi.org/10.69850/rimi.vi3.307}

\bibitem{bhatt2020evaluating}
Bhatt U, Weller A, Moura JMF. Evaluating and aggregating feature-based model
explanations. In: Proceedings of the 29th International Joint Conference on
Artificial Intelligence (IJCAI); 2020. p. 3016--3022.
\url{https://doi.org/10.24963/ijcai.2020/417}

\bibitem{samek2017evaluating}
Samek W, Binder A, Montavon G, Lapuschkin S, M\"{u}ller KR. Evaluating the
visualization of what a deep neural network has learned. IEEE Transactions on
Neural Networks and Learning Systems. 2017;28(11):2660--2673.
\url{https://doi.org/10.1109/TNNLS.2016.2599820}

\bibitem{hooker2019benchmark}
Hooker S, Erhan D, Kindermans PJ, Kim B. A benchmark for interpretability
methods in deep neural networks. In: Advances in Neural Information
Processing Systems 32; 2019. p. 9737--9748.

\bibitem{alvarezmelis2018robustness}
Alvarez-Melis D, Jaakkola TS. On the robustness of interpretability methods.
In: ICML Workshop on Human Interpretability in Machine Learning (WHI 2018);
2018. \url{https://arxiv.org/abs/1806.08049}

\bibitem{hedstrom2023quantus}
Hedstrom A, Weber L, Bareeva D, Krakowczyk D, Motzkus F, Samek W, et al.
Quantus: An explainable AI toolkit for responsible evaluation of neural
network explanations and beyond. Journal of Machine Learning Research.
2023;24:1--11. \url{https://jmlr.org/papers/v24/22-0142.html}

\bibitem{agarwal2022openxai}
Agarwal C, Ley D, Krishna S, Saxena E, Pawelczyk M, Johnson N, et al.
OpenXAI: Towards a transparent evaluation of model explanations. In: 36th
Conference on Neural Information Processing Systems, Datasets and Benchmarks
Track; 2022. \url{https://openreview.net/forum?id=MU2495w47rz}

\bibitem{adebayo2018sanity}
Adebayo J, Gilmer J, Muelly M, Goodfellow I, Hardt M, Kim B. Sanity checks
for saliency maps. In: Advances in Neural Information Processing Systems 31;
2018. p. 9525--9536.

\bibitem{rong2022consistent}
Rong Y, Leemann T, Borisov V, Kasneci G, Kasneci E. A consistent and
efficient evaluation strategy for attribution methods. In: Proceedings of the
39th International Conference on Machine Learning; 2022. p. 18770--18795.
\url{https://proceedings.mlr.press/v162/rong22a.html}

\bibitem{slack2020fooling}
Slack D, Hilgard S, Jia E, Singh S, Lakkaraju H. Fooling LIME and SHAP:
Adversarial attacks on post hoc explanation methods. In: Proceedings of the
AAAI/ACM Conference on AI, Ethics, and Society; 2020. p. 180--186.
\url{https://doi.org/10.1145/3375627.3375830}

\bibitem{kumar2020problems}
Kumar IE, Venkatasubramanian S, Scheidegger C, Friedler SA. Problems with
Shapley-value-based explanations as feature importance measures. In:
Proceedings of the 37th International Conference on Machine Learning; 2020.
p. 5491--5500. \url{https://arxiv.org/abs/2002.11097}

\bibitem{kohavi1996scaling}
Kohavi R. Scaling up the accuracy of Naive-Bayes classifiers: A
decision-tree hybrid. In: Proceedings of the 2nd International Conference on
Knowledge Discovery and Data Mining; 1996. p. 202--207.

\bibitem{becker1996adult}
Becker B, Kohavi R. Adult [dataset]. UCI Machine Learning Repository; 1996.
\url{https://doi.org/10.24432/C5XW20}

\bibitem{breiman2001random}
Breiman L. Random forests. Machine Learning. 2001;45(1):5--32.
\url{https://doi.org/10.1023/A:1010933404324}

\bibitem{chen2016xgboost}
Chen T, Guestrin C. XGBoost: A scalable tree boosting system. In: Proceedings
of the 22nd ACM SIGKDD International Conference on Knowledge Discovery and
Data Mining; 2016. p. 785--794.
\url{https://doi.org/10.1145/2939672.2939785}

\bibitem{demsar2006statistical}
Dem\v{s}ar J. Statistical comparisons of classifiers over multiple data
sets. Journal of Machine Learning Research. 2006;7:1--30.
\url{https://jmlr.org/papers/v7/demsar06a.html}

\bibitem{wilcoxon1945individual}
Wilcoxon F. Individual comparisons by ranking methods. Biometrics Bulletin.
1945;1(6):80--83. \url{https://doi.org/10.2307/3001968}

\bibitem{lakens2013calculating}
Lakens D. Calculating and reporting effect sizes to facilitate cumulative
science: A practical primer for t-tests and ANOVAs. Frontiers in Psychology.
2013;4:863. \url{https://doi.org/10.3389/fpsyg.2013.00863}

\bibitem{holm1979simple}
Holm S. A simple sequentially rejective multiple test procedure.
Scandinavian Journal of Statistics. 1979;6(2):65--70.

\bibitem{ribeiro2018anchors}
Ribeiro MT, Singh S, Guestrin C. Anchors: High-precision model-agnostic
explanations. In: Proceedings of the 32nd AAAI Conference on Artificial
Intelligence; 2018. p. 1527--1535.
\url{https://doi.org/10.1609/aaai.v32i1.11491}

\bibitem{mothilal2020dice}
Mothilal RK, Sharma A, Tan C. Explaining machine learning classifiers through
diverse counterfactual explanations. In: Proceedings of the 2020 Conference
on Fairness, Accountability, and Transparency; 2020. p. 607--617.
\url{https://doi.org/10.1145/3351095.3372850}

\bibitem{dua2019uci}
Dua D, Graff C. UCI Machine Learning Repository. Irvine (CA): University of
California, School of Information and Computer Science; 2019.
\url{http://archive.ics.uci.edu/ml}

\bibitem{wolberg1993breast}
Wolberg W, Mangasarian O, Street N, Street W. Breast cancer Wisconsin
(diagnostic) [dataset]. UCI Machine Learning Repository; 1993.
\url{https://doi.org/10.24432/C5DW2B}

\bibitem{hofmann1994statlog}
Hofmann H. Statlog (German credit data) [dataset]. UCI Machine Learning
Repository; 1994. \url{https://doi.org/10.24432/C5NC77}

\bibitem{haufe2026formalization}
Haufe S, Wilming R, Clark B, Zhumagambetov R, Boubekki A, Martin J, et al.
Explainable AI needs formalization. npj Artificial Intelligence. 2026;2:42.
\url{https://doi.org/10.1038/s44387-026-00095-1}

\bibitem{doshivelez2017rigorous}
Doshi-Velez F, Kim B. Towards a rigorous science of interpretable machine
learning [Preprint]. arXiv:1702.08608; 2017.
\url{https://arxiv.org/abs/1702.08608}

\bibitem{wilming2022scrutinizing}
Wilming R, Budding C, M\"{u}ller KR, Haufe S. Scrutinizing XAI using linear
ground-truth data with suppressor variables. Machine Learning.
2022;111(5):1903--1923. \url{https://doi.org/10.1007/s10994-022-06167-y}

\bibitem{wilming2023theoretical}
Wilming R, Kieslich L, Clark B, Haufe S. Theoretical behavior of XAI methods
in the presence of suppressor variables. In: Proceedings of the 40th
International Conference on Machine Learning; 2023. p. 37091--37107.
\url{https://proceedings.mlr.press/v202/wilming23a.html}

\bibitem{jacovi2020faithfulness}
Jacovi A, Goldberg Y. Towards faithfully interpretable NLP systems: How
should we define and evaluate faithfulness? In: Proceedings of the 58th
Annual Meeting of the Association for Computational Linguistics; 2020.
p. 4198--4205. \url{https://doi.org/10.18653/v1/2020.acl-main.386}

% --- Added in revision 1 (P-11). WP8 verifies every entry and DOI before submission.

\bibitem{bates2015lme4}
Bates D, M\"{a}chler M, Bolker B, Walker S. Fitting linear mixed-effects
models using lme4. Journal of Statistical Software. 2015;67(1):1--48.
\url{https://doi.org/10.18637/jss.v067.i01}

\bibitem{field2007bootstrap}
Field CA, Welsh AH. Bootstrapping clustered data. Journal of the Royal
Statistical Society: Series B. 2007;69(3):369--390.
\url{https://doi.org/10.1111/j.1467-9868.2007.00593.x}

\bibitem{hodges1963estimates}
Hodges JL, Lehmann EL. Estimates of location based on rank tests. The Annals
of Mathematical Statistics. 1963;34(2):598--611.
\url{https://doi.org/10.1214/aoms/1177704172}

\bibitem{kerby2014simple}
Kerby DS. The simple difference formula: An approach to teaching
nonparametric correlation. Comprehensive Psychology. 2014;3:11.IT.3.1.
\url{https://doi.org/10.2466/11.IT.3.1}

\end{thebibliography}

\end{document}
